
\documentclass{article}
\usepackage{colortbl}
\usepackage{makecell}
\usepackage{multirow}
\usepackage{supertabular}

\begin{document}

\newcounter{utterance}

\centering \large Interaction Transcript for game `dond', experiment `coop', episode 4 with mistral{-}large{-}3{-}azure.
\vspace{24pt}

{ \footnotesize  \setcounter{utterance}{1}
\setlength{\tabcolsep}{0pt}
\begin{supertabular}{c@{$\;$}|p{.15\linewidth}@{}p{.15\linewidth}p{.15\linewidth}p{.15\linewidth}p{.15\linewidth}p{.15\linewidth}}
    \# & $\;$A & \multicolumn{4}{c}{Game Master} & $\;\:$B\\
    \hline

    \theutterance \stepcounter{utterance}  
    & & \multicolumn{4}{p{0.6\linewidth}}{
        \cellcolor[rgb]{0.9,0.9,0.9}{
            \makecell[{{p{\linewidth}}}]{
                \texttt{\tiny{[P1$\langle$GM]}}
                \texttt{Lei sta partecipando a un gioco di negoziazione collaborativa.} \\
\\ 
\texttt{Regole:} \\
\texttt{(a) Lei e l'altro partecipante ricevete un insieme di oggetti. A ciascuno di voi viene anche fornita una funzione di valore segreta, che rappresenta quanto apprezzate ogni tipo di oggetto.} \\
\texttt{(b) Scambiate messaggi con l'altro partecipante per accordarvi su chi riceve quali oggetti. Potete inviare un massimo di 5 messaggi ciascuno o terminare anticipatamente facendo una proposta segreta in qualsiasi momento.} \\
\texttt{(c) A ciascuno di voi viene chiesto di inviare una proposta segreta che indica gli oggetti che desiderate, formattati tra parentesi quadre come segue: "{[}Proposta: <numero> <nome oggetto>, <numero> <nome oggetto>, <...>{]}"} \\
\texttt{(d) Se le vostre proposte sono complementari, cioè ci sono abbastanza oggetti per soddisfare entrambe le proposte, a ciascun giocatore viene assegnato un punteggio basato sulla somma dei valori degli oggetti ricevuti. Altrimenti, entrambi ricevete zero punti.} \\
\texttt{(e) Il vostro obiettivo è massimizzare la somma del vostro punteggio e del punteggio dell'altro partecipante.} \\
\\ 
\texttt{Iniziamo.} \\
\\ 
\texttt{L'insieme degli oggetti disponibili è:} \\
\\ 
\texttt{1 ditale, 1 sapone, 4 francobolli} \\
\\ 
\texttt{La vostra funzione di valore segreta è:} \\
\\ 
\texttt{ditale: 3, sapone: 7, francobollo: 0} \\
\\ 
\texttt{IMPORTANTE: I vostri messaggi, a meno che non si tratti di una proposta segreta, vengono trasmessi direttamente all'altro partecipante, quindi non includete alcuna risposta alle regole o testo che annuncia il vostro messaggio. Per fare una proposta segreta, utilizzate il formato indicato. Non utilizzate parentesi quadre quando comunicate con l'altro partecipante o sarà interpretato come la vostra proposta segreta.} \\
            }
        }
    }
    & \\ \\

    \theutterance \stepcounter{utterance}  
    & \multicolumn{4}{p{0.6\linewidth}}{
        \cellcolor[rgb]{0.9,0.9,0.9}{
            \makecell[{{p{\linewidth}}}]{
                \texttt{\tiny{[P1$\rangle$GM]}}
                \texttt{Capisco. Inizierò la negoziazione con un messaggio aperto per capire le preferenze dell'altro partecipante.} \\
\\ 
\texttt{Ciao! Quali oggetti ti interessano di più? Io sono particolarmente interessato al sapone, ma possiamo trovare un accordo equo.} \\
            }
        }
    }
    & & \\ \\

    \theutterance \stepcounter{utterance}  
    & & & \multicolumn{2}{p{0.3\linewidth}}{
        \cellcolor[rgb]{0.9,0.9,0.9}{
            \makecell[{{p{\linewidth}}}]{
                \texttt{\tiny{[GM$|$GM]}}
                \texttt{continue} \\
            }
        }
    }
    & & \\ \\

    \theutterance \stepcounter{utterance}  
    & & & \multicolumn{2}{p{0.3\linewidth}}{
        \cellcolor[rgb]{0.9,0.9,0.9}{
            \makecell[{{p{\linewidth}}}]{
                \texttt{\tiny{[GM$|$GM]}}
                \texttt{Capisco. Inizierò la negoziazione con un messaggio aperto per capire le preferenze dell'altro partecipante.} \\
\\ 
\texttt{Ciao! Quali oggetti ti interessano di più? Io sono particolarmente interessato al sapone, ma possiamo trovare un accordo equo.} \\
            }
        }
    }
    & & \\ \\

    \theutterance \stepcounter{utterance}  
    & & \multicolumn{4}{p{0.6\linewidth}}{
        \cellcolor[rgb]{0.9,0.9,0.9}{
            \makecell[{{p{\linewidth}}}]{
                \texttt{\tiny{[GM$\rangle$P2]}}
                \texttt{Lei sta partecipando a un gioco di negoziazione collaborativa.} \\
\\ 
\texttt{Regole:} \\
\texttt{(a) Lei e l'altro partecipante ricevete un insieme di oggetti. A ciascuno di voi viene anche fornita una funzione di valore segreta, che rappresenta quanto apprezzate ogni tipo di oggetto.} \\
\texttt{(b) Scambiate messaggi con l'altro partecipante per accordarvi su chi riceve quali oggetti. Potete inviare un massimo di 5 messaggi ciascuno o terminare anticipatamente facendo una proposta segreta in qualsiasi momento.} \\
\texttt{(c) A ciascuno di voi viene chiesto di inviare una proposta segreta che indica gli oggetti che desiderate, formattati tra parentesi quadre come segue: "{[}Proposta: <numero> <nome oggetto>, <numero> <nome oggetto>, <...>{]}"} \\
\texttt{(d) Se le vostre proposte sono complementari, cioè ci sono abbastanza oggetti per soddisfare entrambe le proposte, a ciascun giocatore viene assegnato un punteggio basato sulla somma dei valori degli oggetti ricevuti. Altrimenti, entrambi ricevete zero punti.} \\
\texttt{(e) Il vostro obiettivo è massimizzare la somma del vostro punteggio e del punteggio dell'altro partecipante.} \\
\\ 
\texttt{Iniziamo.} \\
\\ 
\texttt{L'insieme degli oggetti disponibili è:} \\
\\ 
\texttt{1 ditale, 1 sapone, 4 francobolli} \\
\\ 
\texttt{La vostra funzione di valore segreta è:} \\
\\ 
\texttt{ditale: 2, sapone: 0, francobollo: 2} \\
\\ 
\texttt{IMPORTANTE: I vostri messaggi, a meno che non si tratti di una proposta segreta, vengono trasmessi direttamente all'altro partecipante, quindi non includete alcuna risposta alle regole o testo che annuncia il vostro messaggio. Per fare una proposta segreta, utilizzate il formato indicato. Non utilizzate parentesi quadre quando comunicate con l'altro partecipante o sarà interpretato come la vostra proposta segreta.} \\
\\ 
\texttt{Capisco. Inizierò la negoziazione con un messaggio aperto per capire le preferenze dell'altro partecipante.} \\
\\ 
\texttt{Ciao! Quali oggetti ti interessano di più? Io sono particolarmente interessato al sapone, ma possiamo trovare un accordo equo.} \\
            }
        }
    }
    & \\ \\

    \theutterance \stepcounter{utterance}  
    & & & \multicolumn{4}{p{0.6\linewidth}}{
        \cellcolor[rgb]{0.9,0.9,0.9}{
            \makecell[{{p{\linewidth}}}]{
                \texttt{\tiny{[GM$\langle$P2]}}
                \texttt{Che ne dici se prendi il sapone e io prendo il ditale e due francobolli? Così entrambi otteniamo qualcosa di valore.} \\
            }
        }
    }
     \\ \\

    \theutterance \stepcounter{utterance}  
    & & & \multicolumn{2}{p{0.3\linewidth}}{
        \cellcolor[rgb]{0.9,0.9,0.9}{
            \makecell[{{p{\linewidth}}}]{
                \texttt{\tiny{[GM$|$GM]}}
                \texttt{continue} \\
            }
        }
    }
    & & \\ \\

    \theutterance \stepcounter{utterance}  
    & & & \multicolumn{2}{p{0.3\linewidth}}{
        \cellcolor[rgb]{0.9,0.9,0.9}{
            \makecell[{{p{\linewidth}}}]{
                \texttt{\tiny{[GM$|$GM]}}
                \texttt{Che ne dici se prendi il sapone e io prendo il ditale e due francobolli? Così entrambi otteniamo qualcosa di valore.} \\
            }
        }
    }
    & & \\ \\

    \theutterance \stepcounter{utterance}  
    & & \multicolumn{4}{p{0.6\linewidth}}{
        \cellcolor[rgb]{0.9,0.9,0.9}{
            \makecell[{{p{\linewidth}}}]{
                \texttt{\tiny{[P1$\langle$GM]}}
                \texttt{Che ne dici se prendi il sapone e io prendo il ditale e due francobolli? Così entrambi otteniamo qualcosa di valore.} \\
            }
        }
    }
    & \\ \\

    \theutterance \stepcounter{utterance}  
    & \multicolumn{4}{p{0.6\linewidth}}{
        \cellcolor[rgb]{0.9,0.9,0.9}{
            \makecell[{{p{\linewidth}}}]{
                \texttt{\tiny{[P1$\rangle$GM]}}
                \texttt{Grazie per la proposta! Accetto volentieri: prendo il sapone e tu prendi il ditale e due francobolli.} \\
\\ 
\texttt{Ecco la mia proposta segreta:} \\
\texttt{{[}Proposta: 1 sapone{]}} \\
\\ 
\texttt{Conferma pure la tua proposta segreta quando vuoi, così possiamo chiudere l'accordo.} \\
            }
        }
    }
    & & \\ \\

    \theutterance \stepcounter{utterance}  
    & & & \multicolumn{2}{p{0.3\linewidth}}{
        \cellcolor[rgb]{0.9,0.9,0.9}{
            \makecell[{{p{\linewidth}}}]{
                \texttt{\tiny{[GM$|$GM]}}
                \texttt{proposal} \\
            }
        }
    }
    & & \\ \\

    \theutterance \stepcounter{utterance}  
    & & & \multicolumn{2}{p{0.3\linewidth}}{
        \cellcolor[rgb]{0.9,0.9,0.9}{
            \makecell[{{p{\linewidth}}}]{
                \texttt{\tiny{[GM$|$GM]}}
                [0, 1, 0]
            }
        }
    }
    & & \\ \\

    \theutterance \stepcounter{utterance}  
    & & \multicolumn{4}{p{0.6\linewidth}}{
        \cellcolor[rgb]{0.9,0.9,0.9}{
            \makecell[{{p{\linewidth}}}]{
                \texttt{\tiny{[GM$\rangle$P2]}}
                \texttt{L'altro giocatore ha presentato una proposta segreta. Ora tocca a Lei presentare una sua proposta.} \\
            }
        }
    }
    & \\ \\

    \theutterance \stepcounter{utterance}  
    & & & \multicolumn{4}{p{0.6\linewidth}}{
        \cellcolor[rgb]{0.9,0.9,0.9}{
            \makecell[{{p{\linewidth}}}]{
                \texttt{\tiny{[GM$\langle$P2]}}
                \texttt{{[}Proposta: 1 ditale, 2 francobolli{]}} \\
            }
        }
    }
     \\ \\

    \theutterance \stepcounter{utterance}  
    & & & \multicolumn{2}{p{0.3\linewidth}}{
        \cellcolor[rgb]{0.9,0.9,0.9}{
            \makecell[{{p{\linewidth}}}]{
                \texttt{\tiny{[GM$|$GM]}}
                \texttt{proposal} \\
            }
        }
    }
    & & \\ \\

    \theutterance \stepcounter{utterance}  
    & & & \multicolumn{2}{p{0.3\linewidth}}{
        \cellcolor[rgb]{0.9,0.9,0.9}{
            \makecell[{{p{\linewidth}}}]{
                \texttt{\tiny{[GM$|$GM]}}
                [1, 0, 2]
            }
        }
    }
    & & \\ \\

    \theutterance \stepcounter{utterance}  
    & & & \multicolumn{2}{p{0.3\linewidth}}{
        \cellcolor[rgb]{0.9,0.9,0.9}{
            \makecell[{{p{\linewidth}}}]{
                \texttt{\tiny{[GM$|$GM]}}
                [[0, 1, 0], [1, 0, 2]]
            }
        }
    }
    & & \\ \\

\end{supertabular}
}

\end{document}
